\begin{figure}[htbp]
\centering
\begin{tikzpicture}[>=stealth]
\draw[arkivblaa] (0,0) circle(1);
\draw[->,black!55] (-1.3,0)--(1.3,0);
\draw[->,black!55] (0,-1.3)--(0,1.3);
\draw[->] (1.7,0)--(3.1,0) node[midway,above] {$A$};
\begin{scope}[shift={(5.5,0)},rotate=25]
\draw[arkivblaa] (0,0) ellipse[x radius=1.8,y radius=.65];
\draw[->,black!65] (0,0)--(1.8,0) node[right] {$\sigma_1u_1$};
\draw[->,black!65] (0,0)--(0,.65) node[above] {$\sigma_2u_2$};
\end{scope}
\end{tikzpicture}
\caption{SVD i et reelt todimensionalt rum: Enhedscirklen afbildes på en ellipse. $V^{\mathsf T}$ vælger inputretninger, $\Sigma$ skalerer med singularværdierne, og $U$ orienterer resultatet. Figuren viser $\sigma_1=1{,}8$ og $\sigma_2=0{,}65$ med samme dimensionsløse målestok.}
\label{fig:b1-svd}
\end{figure}
